% References: https://en.cppreference.com/w/cpp/language/operators.html
% https://eel.is/c++draft/associative.reqmts
% https://eel.is/c++draft/priority.queue
\subsection{重载运算符与自定义比较}
\label{stl:operators}
\subsubsection*{原理}
给结构体定义 \InlineCode{operator<}，即可让默认排序、有序集合和堆知道怎样比较对象。成员写法中，当前对象是左操作数，参数是右操作数；末尾的 \InlineCode{const} 表示比较不修改当前对象，参数用常量引用避免复制。

比较必须形成严格弱序，不能把小于写成小于等于。\InlineCode{set/map} 判断两个键等价的条件是两方向比较均为假，\textbf{不依赖} \InlineCode{operator==}；若同权值但不同编号的点应分别保留，就把编号也作为次关键字。
\subsubsection*{使用教程}
以下按 \InlineCode{x}、再按 \InlineCode{id} 升序比较，可直接用于排序、集合和映射。\InlineCode{operator>} 复用反向小于关系，便于构造 \InlineCode{greater<Node>} 小根堆；\InlineCode{operator==} 判断字段相等，供 \InlineCode{find/unique} 等使用。C++17 不会仅凭小于自动生成其余比较运算符。
\begin{lstlisting}[style=mtycpp,numbers=none]
struct Node{
	long long x;int id;
	bool operator<(const Node &b)const{
		if(x!=b.x) return x<b.x;
		return id<b.id;
	}
	bool operator>(const Node &b)const{return b<*this;}
	bool operator==(const Node &b)const{
		return x==b.x&&id==b.id;
	}
};
vector<Node> a={{3,2},{1,5},{3,1}};
sort(a.begin(),a.end());
set<Node> st(a.begin(),a.end());
map<Node,int> cnt;
cnt[{3,1}]++;
priority_queue<Node> mx;
priority_queue<Node,vector<Node>,greater<Node>> mn;
for(auto p:a){mx.push(p);mn.push(p);}
struct Cmp{
	bool operator()(const Node &a,const Node &b)const{
		return b<a;
	}
};
set<Node,Cmp> desc(a.begin(),a.end());
priority_queue<Node,vector<Node>,Cmp> heap;
heap.push({3,1});
\end{lstlisting}
自然的小于关系下，\InlineCode{sort} 和 \InlineCode{set/map} 按升序排列，默认 \InlineCode{priority_queue} 却取最大者；改用 \InlineCode{greater} 或上述 \InlineCode{Cmp} 后，堆顶取最小者。同一个 \InlineCode{Cmp} 用在集合上则表示降序。\InlineCode{operator()} 让比较器对象可像函数一样接收两个参数，适合只改变某个容器的规则。

只为小根堆把 \InlineCode{operator<} 写反也可行，但该类型再放进默认 \InlineCode{set} 或 \InlineCode{sort} 时顺序也会反转。需要多种顺序时，保留自然的小于并单独定义比较器更直观。\InlineCode{unordered_map/set} 还需要相应的哈希函数与相等判断，仅重载小于不能作为哈希键直接使用。
